\chapter{Group Evaluation}
This chapter presents our subjective views of the project.

\section{Introduction}
In the course of the project, the group has worked well as both a team and as
individuals, drawing from each other strengths and expertise. We have created
something that could very well either live on in service of the information
security community or inspire others to further improve or recreate.

\section{Organization}
The group was given an interesting task of to look at an existing software
prototype, and build a completely new system with added functionality and better
performance. In addition, \snortmanager{}, which is a very similar project, had
already been conducted at \GUC{}. It was therefore important not only to create
a solid system, but to develop a software that is clearly different than any
existing solutions. This resulted in a thourough design-phase that laid a solid
foundation for the development.

While we had a designated group leader, we never really got in a situation where
we needed one. Most of our decisions or problems were solved as a group through
discussion and a democratic approach and that worked well for our group. It
seems like we succeeded when it came to our project plan, which we have more or
less followed throughout the entire project. The development cycle lengths and
intermittent meetings have allowed us to be very flexible and agile in our
development, as well as being more or less according to schedule. Spending some
time planning everything before we started programming gave us a good head start
and even saved us some time writing this report as we had a lot of its
components already made. When we completed the initial planning phase, we held a
``project decision meeting'' with Avdeling BKI where we introduced the system
models we had made so that Avdeling BKI could decide what they wanted or not. We
felt this was a good idea, as we could then get some valuable feedback before we
started programming.

Our approach to the development cycles was mostly as described in section
\ref{sec:dev_method} we conducted synchronisational meetings once or twice a
week to chart out the course, and keep up to date with what the others were
doing. But for the most part, because we had divided the software into modules, we worked
individually. This worked out well, as the group members had different schedules
and needed the flexibility. For the thesis, we set a good month and a half of
our project time for it and we think it paid off, especially as we expanded it to
include the performance testing part. We also divided the writing part into
three draft cycles, so we could get valuable feedback from our supervisor and
Avdeling BKI.

\section{Work distribution}
For the most part our work distribution has been a ``free for all'', meaning
the group members picked freely from what they wanted to work on during the
development cycle, within the objectives we set for that cycle. 
This worked out fine for the most part, but we feel we probably could have been
slightly more aggressive with how many features we finished within the
cycles. This left us with a bigger workload towards the end of the development
as opposed to the beginning.

\section{What could have been done differently}
In hindsight, we probably could have conducted our performance testing earlier
in the project, so that we could have more time to implement optimizations. We
should have also had a more strict approach to what we were supposed to work on
for the different development cycles, planning that out earlier as well.

\section{Subjective views}

\subsection{Thomas Nyheim}
For me, this project has had a fair bit of personal interest in its success as I
am the one that is going to be using this while working for Avdeling BKI. That
meant that the better this project turned out, the better for me it would be
afterwards. 

I feel we have created something that works well enough that it can be taken
further and improved upon still. We've also conducted a small scientific test
case which has been a learning experience, as it is not something often done at
Bachelor levels at \GUC. There has also been a learning curve when it comes to
Python and Django, which I can take with me in the future and apply to other
endevours.

\subsection{Eigil Obrestad}
***Skriv om folelser og slikt***
\subsubsection{Eirik Skogstad}
***Skriv om folelser og slikt***